\documentclass[11pt,a4paper]{article}
\usepackage[utf8]{inputenc}
\usepackage[T1]{fontenc}
\usepackage{lmodern}
\usepackage[margin=1in]{geometry}
\usepackage{amsmath,amssymb,amsthm}
\usepackage{booktabs}
\usepackage{longtable}
\usepackage{array}
\usepackage{microtype}
\usepackage[hidelinks]{hyperref}

\newtheorem{theorem}{Theorem}
\newtheorem{proposition}[theorem]{Proposition}
\newtheorem{lemma}[theorem]{Lemma}
\newtheorem{corollary}[theorem]{Corollary}
\theoremstyle{definition}
\newtheorem{observation}[theorem]{Observation}
\theoremstyle{remark}
\newtheorem{remark}[theorem]{Remark}

\newcommand{\tier}[1]{{\small\textsf{[#1]}}}
\newcommand{\T}{T_R}
\newcommand{\lmin}{\lambda_{\min}}
\newcommand{\lstar}{\lambda^{*}}
\newcommand{\Fq}{\mathbb{F}_q}
\newcommand{\src}[1]{{\footnotesize\texttt{#1}}}

\title{\bfseries A $\delta^2$--$L^3$ Sensitivity Law for Finite-Window Weil
Positivity, with a Function-Field Control of the Connes--van Suijlekom Pipeline}

\author{Elias De Jes\'us\\[2pt]
{\small Independent Researcher}\\[1pt]
{\small ORCID \href{https://orcid.org/0009-0007-0190-9143}{0009-0007-0190-9143}}\\[1pt]
{\small \texttt{dejesuselias10@gmail.com}}}

\date{October 2026 --- Technical Note\\[2pt]
{\small\href{https://doi.org/10.5281/zenodo.23195177}{doi:10.5281/zenodo.23195177}}}

\begin{document}
\maketitle

\begin{abstract}
\noindent
The Connes--van Suijlekom and Connes--Consani--Moscovici programs reduce the Riemann
hypothesis to spectral properties of a quadratic form restricted to test functions
supported in a finite window. We ask how sensitive that form is to a zero moved off the
critical line, and we calibrate the answer in a setting where the truth is known: curves
over finite fields, where the window form is exactly the Frobenius Gram matrix of
Hallouin--Perret and the Riemann hypothesis is Weil's theorem. Our main result is a
sensitivity law. In the function field we prove that moving a reciprocal real pair
$\{\rho,\rho^{-1}\}$ off the critical circle gives
$\lmin(\T)=-\binom{R+2}{3}a^{2}+O(a^{4})$ with $a=\log\rho$, the coefficient being twice
the centred second moment of the window. On the zeta side, in the even sector, we prove an
bound $\lmin(Q_\delta)\ge\lstar(L)-c(L)\delta^{2}-O(\delta^{4})$ with
$c(L)=8L^{3}(\tfrac13+\tfrac1{\sqrt5})$, stated \emph{relative} to a lower bound $\lstar(L)$ for
the unperturbed form and unconditional where that positivity is certified ($L\le0.8$), and
a sharper conditional bound with the \emph{exact} constant $4K(L,\gamma)$, where
$K(L,\gamma)=\int_{-L}^{L}u^{2}\sin^{2}(\gamma u)\,du$ is attained by
$f\propto u\sin(\gamma u)$; numerically the re-optimised minimiser reaches up to $96\%$ of
it. Under an explicit dictionary the two laws are the same expression,
$(\text{window length})^{3}/6$. We also run the full Connes--van Suijlekom minimiser
pipeline on deliberately planted off-circle spectra, which separates the steps that carry
arithmetic content from those that do not.
\textbf{Nothing here proves or advances the Riemann hypothesis for $\zeta$.}
\end{abstract}

\section*{Non-claims}
\addcontentsline{toc}{section}{Non-claims}

Stated first, not in an appendix.

\begin{enumerate}\itemsep2pt
\item \textbf{No proof or advance of RH for $\zeta$.} Part~I runs where the Riemann
hypothesis is a theorem (Weil, curves over finite fields); Part~II uses fabricated spectra.
\item \textbf{No new spectral architecture.} The window form, its Toeplitz structure and
its Carath\'eodory--Fej\'er theory are due to Hallouin--Perret~\cite{HP19} and
Connes--van Suijlekom~\cite{CvS}; the finite-window Weil form itself goes back to
Yoshida~\cite{Yoshida} and Bombieri~\cite{Bombieri}.
\item \textbf{Planted spectra are fabricated diagnostics}, not $L$-functions of anything.
\item \textbf{Numerical detections are tier T2 unless certified}, and
\emph{non-detections certify nothing}: a finite basis that finds no negative direction is
not evidence of positivity. This asymmetry is the subject of \S\ref{sec:onset}.
\item \textbf{$\lstar(0.8)\le 2.2702\times10^{-17}$ is a variational upper bound} from a
finite basis, not a certified enclosure. We prove no rigorous lower bound \emph{of our own};
Zhu~\cite{Zhu} certifies one, and Theorem~\ref{thm:zeta} uses it to convert its relative
estimate into an unconditional lower bound for $L\le0.8$. Theorem~\ref{thm:zetasharp} does
\emph{not} use it; it rests on its own hypothesis.
\item \textbf{\S\ref{sec:zeta} reproduces published computations} (Groskin~\cite{Groskin1})
at smaller scale.
\item \textbf{Block F is claimed only for the function field.} On $\zeta$ windows the
indefiniteness of $A$ and $M$ depends on where $-\log\pi$ is placed; see
\S\ref{sec:recorder} and Retraction~15.
\item \textbf{Not attempted:} Davenport--Heilbronn; any Lean formalisation.
\end{enumerate}

\section*{Epistemic tiers}
\addcontentsline{toc}{section}{Epistemic tiers}

\begin{center}
\begin{tabular}{@{}lll@{}}
\toprule
Tier & Meaning & Discharged by \\
\midrule
T1 & exact or proved & a proof, or exact arithmetic over $\mathbb{Z}$ or $\mathbb{Q}$ \\
T2 & numerical & stated working precision $+$ a convergence or noise-floor check \\
T3 & interpretation & argument from T1/T2 evidence, labelled as such \\
\bottomrule
\end{tabular}
\end{center}

\noindent Every numbered statement below carries a tier, the script that produced it and
the data file that records it. All numerics use \texttt{mpmath} at the stated precision;
\texttt{float64} is used only as a deliberate contrast (\S\ref{sec:traps}).

\section{The pipeline and the question}

Let $X/\Fq$ be a smooth projective absolutely irreducible curve of genus $g$, and let
$N_n=\#X(\mathbb{F}_{q^n})$. Hallouin and Perret~\cite[Prop.~5]{HP19} show that the
normalised classes $\gamma^{k}$ of the iterated Frobenius graphs have Gram matrix
\begin{equation}\label{eq:HPgram}
\operatorname{Gram}(\gamma^0,\dots,\gamma^R)=\bigl(t(|i-j|)\bigr)_{0\le i,j\le R},
\qquad
t(0)=2g,\quad t(n)=\frac{(q^{n}+1)-N_n}{q^{n/2}},
\end{equation}
a symmetric Toeplitz matrix, which we write $\T$. We adopt~\eqref{eq:HPgram} as the
definition of the window form; it is \emph{their} matrix, entry for entry.

The Connes--van Suijlekom pipeline~\cite{CvS} passes through four steps: (i) reality of the
zeros of the ground-state polynomial; (ii) the sign of $\lmin$; (iii) simplicity and
evenness of the ground state; (iv) convergence of the normalised minimiser. The question
of this note is \emph{quantitative}: how large a negative eigenvalue does a zero moved a
distance $\delta$ off the critical line actually produce, and how large a window is needed
to see it?

\paragraph{What is claimed as new.} Only two things.
\begin{enumerate}
\item[(N1)] The $\delta^{2}$--$L^{3}$ sensitivity law: Theorem~\ref{thm:ff} (function field),
Theorems~\ref{thm:zeta} and~\ref{thm:zetasharp} ($\zeta$), the measured near-saturation,
and the shared $(\text{window length})^{3}/6$ form.
\item[(N2)] The planted-spectrum control of the CvS minimiser pipeline in the function
field (\S\ref{sec:planted}).
\end{enumerate}
Everything else is foundation or confirmation; \S\ref{sec:prior} says exactly whose.

\section{Foundation (prior art, stated as such)}\label{sec:found}

\begin{proposition}[Hallouin--Perret~{\cite[Prop.~5, \S1.2]{HP19}}]\label{prop:HP5}
\tier{T1, cited} The Gram matrix of the normalised iterated-Frobenius classes is the
Toeplitz matrix~\eqref{eq:HPgram}.
\end{proposition}

\begin{theorem}[Hallouin--Perret~{\cite[Thm.~6, \S1.3]{HP19}}]\label{thm:HP6}
\tier{T1, cited} Let $d$ be the rank of the Frobenius space $F$. Then $d$ is the minimal
integer such that $\operatorname{Gram}(\gamma^0,\dots,\gamma^{d})$ is singular, a generator
of that kernel exists over $\mathbb{Q}$, and the resulting Frobenius eigenvalues are
algebraic integers of modulus $\sqrt q$.
\end{theorem}

\begin{theorem}[Hallouin--Perret~{\cite[Thm.~36(ii), App.~A.2]{HP19}}]\label{thm:HP36}
\tier{T1, cited} If a real symmetric positive semi-definite Toeplitz matrix of size $n+1$
has rank $n$ and kernel vector $(a_0,\dots,a_n)$, then $P=\sum a_jX^j$ has $n$ distinct
roots, all of modulus $1$.
\end{theorem}

Theorem~\ref{thm:HP36} is step~(i) of the pipeline; it is the Carath\'eodory--Fej\'er
corollary that CvS~\cite[Cor.~1.1]{CvS} prove again $C^{*}$-algebraically. Their
Lemma~33 (rank $=$ largest non-zero \emph{leading} minor) and Lemma~34 (the
$(\gamma_i\pm\gamma_{n-i})/\sqrt2$ parity split) are the computational tools we use
throughout. CvS themselves point to~\cite{HP19} for the function-field connection and
write that ``the key difficulty in this context, then, becomes the verification that zero
is indeed the (simple) minimal eigenvalue of $T$''~\cite[Introduction]{CvS}.

\section{The control experiment}\label{sec:planted}

\subsection{Genuine spectra}
\tier{T1 for the arithmetic, T2 for the spectra; \src{src/stage1.py}, \src{data/stage1\_spectra.csv}}
Four curves ($y^2=x^3+x/\mathbb{F}_5$, $y^2=x^5+x^2+1/\mathbb{F}_3$,
$y^2=x^7+x+1/\mathbb{F}_5$, $y^2=x^7+x^3+x+1/\mathbb{F}_3$; $g=1,2,3$) with brute-force
point counts, $L$-polynomials by Newton's identities plus the functional equation, and
exact integer power sums. All $40$ rows have $\lmin\ge0$; $\lmin(R)$ decreases monotonically
(Cauchy interlacing) and vanishes exactly at $R=d$, where the kernel polynomial's roots are
the $e^{i\theta_j}$ --- Theorem~\ref{thm:HP6} in action.

\subsection{Planted spectra}
\tier{T1, exact over $\mathbb{Q}$; \src{src/stage2.py}, \src{src/exact\_inertia.py}}
We replace the Frobenius spectrum by a functional-equation-closed multiset with
$|\beta|\ne1$. Choosing rational $\rho$ and $\cos\varphi$ makes every $t(n)$ rational, so
the inertia is certified \emph{exactly} by the Sylvester--Jacobi leading-minor rule.

\begin{center}
\begin{tabular}{@{}llccc@{}}
\toprule
configuration & planted & first $R$ with $\lmin<0$ & step~(i) holds? & $\lmin$ simple? \\
\midrule
on-line control & $|\beta|=1$ & never (PSD, proved) & always & fails for every $R>d$ \\
real pair & $\rho=1.3$ & $R=1$ & always & never fails \\
quartet & $\rho=1.3,\varphi=0.7$ & $R=2$ & always & never fails \\
two quartets & $\rho=1.3$ & $R=4$ & always & never fails \\
mixed & quartet $+$ pair & $R=4$ & always & never fails \\
\bottomrule
\end{tabular}
\end{center}

\begin{observation}\label{obs:free}\tier{T1 reasoning, T2 evidence}
Step~(i) holds at every applicable row of every planted configuration, including rows where
$\lmin=-642.8$: measured $\max\bigl||z|-1\bigr|\le5.1\times10^{-48}$. Steps~(i) and~(iii)
are therefore free of arithmetic content. This is not a discovery --- CvS state
Theorem~\ref{thm:HP36} for an arbitrary real even distribution, so no arithmetic can enter
step~(i) --- but the planted control makes it concrete, with numbers attached.
\end{observation}

\subsection{Degeneracy beyond the window degree}
\tier{T1 corollary of Thm.~\ref{thm:HP36}(i); \src{src/function\_level.py}}
Since the shift has $P$ as both minimal and characteristic polynomial, $\ker\T$ is the
ideal generated by $P$ truncated at degree $R$:
\begin{equation}
\ker\T=\{P\cdot q:\deg q\le R-d\},\qquad \dim\ker\T=\max(0,R+1-d),
\end{equation}
matched in every computed row. Hence $\lmin=0$ is simple exactly at $R=d$ and never after,
so the CvS simple-ground-state hypothesis fails for every $R>d$; and since $P$ is
self-reciprocal, the kernel splits into even and odd parts of dimensions
$\lceil\frac{m+1}2\rceil,\lfloor\frac{m+1}2\rfloor$ with $m=R-d$, both non-zero for
$m\ge1$. CvS state this in words in their Remark~2.3; we record the dimensions.

\subsection{Step (iv) at the function level}
\tier{T1; \src{src/function\_level.py}}
\begin{proposition}\label{prop:hurwitz}
If $\lmin$ is simple for infinitely many $R$ and the normalised ground-state polynomials
converge locally uniformly to $F\not\equiv0$, then every zero of $F$ lies on $|z|=1$ by
Theorem~\ref{thm:HP36} and Hurwitz's theorem. Hence they cannot converge to a target with
a zero off the circle.
\end{proposition}
\noindent The free witness is exactly the obstruction. Measured: zero \emph{angles}
converge to $\sim10^{-3}$ even for off-line spectra, while the coefficient distance to the
target is exactly $\sqrt2$ (orthogonal) and the normalised $|\hat c_R(\beta)|$ \emph{rises}
to $0.58$--$0.71$ out to $R=36$ instead of falling. This makes concrete the Hurwitz step
that is CvS's own Step~5.

\section{The $\delta^2$--$L^3$ law}\label{sec:law}

\subsection{Function field}

\begin{theorem}\label{thm:ff}\tier{T1; proof in \src{notes/proofs.md}, checked by \src{src/proofs\_check.py}}
Let the normalised spectrum be a single off-circle real pair $\{\rho,\rho^{-1}\}$ and put
$a=\log\rho$. With $u_i=\rho^{i}$, $w_i=\rho^{-i}$ $(0\le i\le R)$,
\[
\T=uw^{\mathsf T}+wu^{\mathsf T},\qquad
\lmin(\T)=(R+1)-\|u\|\,\|w\|,
\]
all other eigenvalues being $0$. Consequently $\lmin<0$ for every $R\ge1$ as soon as
$\rho\ne1$, by strict Cauchy--Schwarz; $\lmin$ is an even function of $a$; and
\[
\boxed{\ \lmin(\T)=-2a^{2}\sum_{i=0}^{R}\Bigl(i-\tfrac R2\Bigr)^{2}+O(a^{4})
=-\binom{R+2}{3}a^{2}+O(a^{4}).\ }
\]
For $R=1$ this is exact: $\lmin=-4\sinh^{2}(a/2)$.
\end{theorem}

The coefficient is twice the \emph{centred} second moment of the window; the centring falls
out of the optimisation, not from a convention. Checks: the factorisation to $6.2\times10^{-61}$,
the eigenvalues to $1.0\times10^{-59}$, the moment identity exactly for $R=1,\dots,199$, and
$-\lmin/a^{2}=\binom{R+2}{3}$ to twelve printed digits for $R=1,2,4,8,16,32$.

\subsection{Zeta side}

Throughout this subsection $f$ is real, \textbf{even}, supported in $[-L,L]$ with
$\|f\|_2=1$, and $F=\hat f$; all statements here are for the even sector.

\paragraph{The baseline, stated explicitly.} An off-line zero forces a quartet --- four
zeros --- whereas a \emph{simple} on-line zero at $\gamma$ contributes a pair. The zeros of
$\zeta$ are expected to be simple, and every computed one is, so a count-preserving
comparison must be made against a \emph{hypothetical double} on-line zero. Writing
$Q_{\mathrm{true}}$ for the actual Weil form, the baseline below is
\begin{equation}\label{eq:baseline}
Q_0 \;=\; Q_{\mathrm{true}} + 2F(\gamma)^{2}.
\end{equation}
Two consequences, both used. First $Q_0\ge Q_{\mathrm{true}}$, so Theorem~\ref{thm:zeta} is
unaffected: a lower bound certified for $Q_{\mathrm{true}}$ is also one for $Q_0$. Second,
Theorem~\ref{thm:zetasharp}'s hypothesis $\widetilde Q:=Q_0-4F(\gamma)^{2}\ge0$ reads
$Q_{\mathrm{true}}-2F(\gamma)^{2}\ge0$, i.e.\ exactly ``every zero \emph{other} than the one
being moved lies on the critical line''. Of the three conventions this is the only one with
$Q_\delta\to Q_0$ as $\delta\to0$, and~\eqref{eq:baseline} is what that continuity costs: it
is bought by shifting the baseline, not by comparing against $\zeta$ itself.

Moving the double zero to $\tfrac12+\delta\pm i\gamma$ changes the form by
\begin{equation}\label{eq:pert}
Q_\delta(f)-Q_0(f)=-4\delta^{2}\bigl(F'(\gamma)^{2}+F(\gamma)F''(\gamma)\bigr)+R_4 .
\end{equation}

\paragraph{The two baselines, the budget, and what $\lstar$ denotes.} Call $f$
\emph{admissible} if it is real, even, $\mathrm{supp}\,f\subset[-L,L]$ and $\|f\|_2=1$. Write
$Q_\zeta$ for the \emph{actual} $\zeta$ Weil form on the window, so that
$Q_0=Q_\zeta+2F(\gamma)^{2}$ by~\eqref{eq:baseline}, and set
\[
b_0(L,\gamma)=\inf_{f\ \mathrm{admissible}}Q_0(f),
\qquad
\lstar(L)\ \le\ \inf_{f\ \mathrm{admissible}}Q_\zeta(f),
\]
so that $b_0$ is the infimum of the \emph{augmented counterfactual} baseline and $\lstar(L)$
denotes any \emph{lower bound} for the \emph{actual} unperturbed form. These are two different
objects and are nowhere identified. Both are infima of a quadratic form on the stated domain;
nothing here assumes the Weil form is a bounded operator on $L^{2}$. Put
\[
B_L(\delta)=c(L)\delta^{2}+\tfrac{16}{3}L^{5}\delta^{4}e^{2L\delta},
\qquad c(L)=8L^{3}\bigl(\tfrac13+\tfrac1{\sqrt5}\bigr)\approx6.244\,L^{3},
\]
the perturbation budget: $4|F'(\gamma)^{2}+F(\gamma)F''(\gamma)|\le c(L)$ by Cauchy--Schwarz
and $|R_4|\le\tfrac{16}{3}L^{5}\delta^{4}e^{2L\delta}$, both proved in \S B.2--B.3 of
\src{notes/proofs.md}.

A lower bound for $Q_\zeta$ on a larger test space is \emph{a fortiori} a lower bound on the
even sector, because an infimum over a larger space cannot exceed one over a subspace; so a
certificate proved on the full admissible space may be quoted here unchanged. Zhu's
certificate is used only with its documented normalisation $Q(f)/\|f\|^{2}$ --- the one used
throughout --- and only on its documented range.

\begin{theorem}[relative perturbation estimate; even sector]\label{thm:zeta}\tier{T1}
For every admissible $f$ and every $\delta\ge0$,
\[
Q_\delta(f)\ \ge\ Q_0(f)-B_L(\delta),
\]
and hence, taking infima over the normalised admissible sector,
\[
\inf Q_\delta\ \ge\ b_0(L,\gamma)-B_L(\delta)
\qquad\text{and}\qquad
\lmin(Q_\delta)\ \ge\ \lstar(L)-B_L(\delta),
\]
the second because $Q_0\ge Q_\zeta$, so any lower bound $\lstar(L)$ for the actual unperturbed
form is also one for $Q_0$.
\end{theorem}

The estimate holds \emph{relative} to the unperturbed lower bound $\lstar(L)$. Replacing this
baseline by zero requires independently established nonnegativity. We use the cited positivity
certificate for $L\le0.8$: Zhu~\cite{Zhu} certifies $\lstar(0.8)\ge8.9\times10^{-18}$ and
$\lstar$ is non-increasing in $L$, so
\[
\lmin(Q_\delta)\ \ge\ 8.9\times10^{-18}-B_L(\delta)
\qquad (L\le0.8).
\]
No such replacement is asserted here for larger windows.

\paragraph{Validity versus usefulness of the budget.} The remainder bound inside
$B_L(\delta)$ is a convergent series tail, not a truncated asymptotic expansion, so
$B_L(\delta)$ is a valid budget at every $\delta$ on its proved domain. What degrades as
$L\delta$ grows is its \emph{informativeness}: the $\delta^{4}$ term overtakes
$c(L)\delta^{2}$, after which $B_L(\delta)$ no longer reflects the quadratic law. Validity and
usefulness are separate questions and only the second fails there.

\begin{theorem}[sharper, conditional; even sector]\label{thm:zetasharp}\tier{T1}
Define the \emph{remaining-zero form}
\[
Q_{\mathrm{rest}}(f)\;=\;Q_0(f)-4F(\gamma)^{2}\;=\;Q_\zeta(f)-2F(\gamma)^{2},
\]
written $\widetilde Q$ below. Suppose every nontrivial zero other than the one being moved lies
on the critical line; where the explicit formula applies, that is a sufficient input for
$\widetilde Q\ge0$. It is not a necessary one, and $\widetilde Q\ge0$ does \emph{not} follow
from positivity of $Q_\zeta$, from positivity of $Q_0$, or from a zero-free strip: each of
those leaves room for $2F(\gamma)^{2}$ to exceed the margin. Completing the square,
$4F^{2}-4\delta^{2}FF''=(2F-\delta^{2}F'')^{2}-\delta^{4}F''^{2}$, moves the $FF''$
allowance into $O(\delta^{4})$ and gives
\[
\lmin(Q_\delta)\ \ge\ -4\,K(L,\gamma)\,\delta^{2}
-\Bigl(\tfrac25+\tfrac{16}{3}e^{2L\delta}\Bigr)L^{5}\delta^{4},
\]
where $K(L,\gamma)$ is the \emph{exact} value of the extremal problem,
\begin{equation}\label{eq:K}
K(L,\gamma)\;=\!\!\sup_{\substack{f\ \mathrm{real,\ even},\ \mathrm{supp}f\subset[-L,L]\\ \|f\|_2=1}}\!\!|F'(\gamma)|^{2}
\;=\;\int_{-L}^{L}\!u^{2}\sin^{2}(\gamma u)\,du,
\end{equation}
attained by $f\propto u\sin(\gamma u)$, with closed form
\[
K(L,\gamma)=\frac{L^{3}}{3}-\left[\frac{L^{2}\sin 2\gamma L}{2\gamma}
+\frac{L\cos 2\gamma L}{2\gamma^{2}}-\frac{\sin 2\gamma L}{4\gamma^{3}}\right].
\]
\end{theorem}

\begin{corollary}[$\gamma$-uniform form]\label{cor:uniform}\tier{T1}
Since $\sin^{2}\le1$ gives $K(L,\gamma)\le\int_{-L}^{L}u^{2}du=2L^{3}/3$ for every $\gamma$,
\[
\lmin(Q_\delta)\ \ge\ -\tfrac{8L^{3}}{3}\,\delta^{2}-O(\delta^{4}),
\]
which is the form to quote when $\gamma$ is not fixed. As $\gamma\to\infty$,
$K(L,\gamma)\to L^{3}/3$, so the uniform constant overshoots the large-ordinate value by
exactly the factor $2$ coming from the average of $\sin^{2}$.
\end{corollary}

\begin{remark}\label{rem:twoconstants}
$4K(L,\gamma)$ is sharp \emph{for the derivative-evaluation step}: by~\eqref{eq:K} the
supremum is attained, so no smaller $\delta^{2}$ coefficient is available in that step. Three
claims must be kept apart. \textbf{(a)} $K(L,\gamma)$ is the exact squared norm of the
derivative-evaluation functional $f\mapsto F'(\gamma)$ on the normalised admissible sector.
\textbf{(b)} $f\propto u\sin(\gamma u)$ attains that norm in the stated space. \textbf{(c)}
Neither implies that the minimiser of the \emph{entire} Weil form attains the same sensitivity
coefficient --- attaining the derivative norm need not minimise $Q_\delta$. So $4K$ is not
claimed to be a proved sharp coefficient for the optimised minimum eigenvalue; the
near-saturation measured in \S\ref{sec:sat} is T2. The constants $8L^{3}/3$ and $4L^{3}/3$
appearing below are respectively the $\gamma$-uniform upper bound and the $\gamma\to\infty$
limit of the same quantity, and ratios reported against them should be read as such.
\end{remark}

\subsection{How close the measured constant comes}\label{sec:sat}
\tier{T2, \src{src/task\_final\_checks.py}, \src{data/stage3\_exact\_constant\_K.csv}}

The planted ordinate is $\gamma_1=14.1347$. Measured $C(L)$ against the exact constant
$4K(L,\gamma_1)$ of Theorem~\ref{thm:zetasharp}, and for reference against the two
constants of Corollary~\ref{cor:uniform}:

\begin{center}
\begin{tabular}{@{}lccccc@{}}
\toprule
$L$ & $0.8$ & $1.0$ & $1.3$ & $1.6$ & $2.0$ \\
\midrule
measured $C(L)$ & $0.3108$ & $0.8530$ & $2.4987$ & $4.9218$ & $10.083$ \\
$4K(L,\gamma_1)$ & $0.7419$ & $1.3427$ & $3.1158$ & $5.1130$ & $10.652$ \\
$C/4K(L,\gamma_1)$ & $0.419$ & $0.635$ & $0.802$ & $\mathbf{0.963}$ & $0.947$ \\[2pt]
$C/(8L^{3}/3)$ \emph{(uniform)} & $0.228$ & $0.320$ & $0.427$ & $0.451$ & $0.473$ \\
$C/(4L^{3}/3)$ \emph{($\gamma\to\infty$)} & $0.455$ & $0.640$ & $0.853$ & $0.901$ & $0.945$ \\
\bottomrule
\end{tabular}
\end{center}

Against the exact constant the re-optimised minimiser reaches $96.3\%$ at $L=1.6$.

\paragraph{How much the $\sin^{2}$ weight matters, measured.} The deviation
$|K(L,\gamma_1)/(L^{3}/3)-1|$ is
\[
17.3\%,\quad 8.68\%,\quad 0.699\%,\quad 6.37\%,\quad 6.38\%,\quad 0.136\%
\]
at $L=0.6,0.8,1.0,1.3,1.6,2.0$. It is \emph{not} monotone: the bracketed terms
in~\eqref{eq:K} oscillate in $L$ with period $\pi/\gamma$, so the correction can be large at
one window and negligible at a slightly larger one. At the window where positivity is
certified, $L=0.8$, it is $8.68\%$; at $L=2$ it happens to be $0.136\%$. This is why the
$\gamma$-uniform constant should not be used as a stand-in for $K$ at a fixed ordinate.

\subsection{One law, two settings: the dictionary}\label{sec:dict}
\tier{T3, checked numerically in \src{src/task\_final\_checks.py}}

The two laws are the same statement under an explicit change of variables. For a curve over
$\Fq$ the explicit-formula variable is $\log$ of a prime power, so the lattice index $n$ and
the $\zeta$ variable $u$ correspond by $u=n\log q$. This fixes everything else:

\begin{center}
\begin{tabular}{@{}lll@{}}
\toprule
 & function field & $\zeta$ \\
\midrule
window variable & $n=0,\dots,R$ & $u\in[-L,L]$ \\
correspondence & \multicolumn{2}{l}{$u=n\log q$, hence window length $R\log q=2L$} \\
off-line depth & $a=\log\rho$ & $\delta=\mathrm{Re}\,\rho-\tfrac12$, with $a=\delta\log q$ \\
Rayleigh quotient & sum over $n$ & integral $du=\log q\,dn$: differ by $\log q$ \\
\bottomrule
\end{tabular}
\end{center}

Substituting $R=2L/\log q$ and $a=\delta\log q$ into Theorem~\ref{thm:ff} and dividing by
$\log q$:
\[
-\binom{R+2}{3}a^{2}\;\approx\;-\frac{R^{3}}{6}a^{2}
=-\frac{1}{6}\Bigl(\frac{2L}{\log q}\Bigr)^{3}(\delta\log q)^{2}
=-\frac{4L^{3}}{3}\,\frac{\delta^{2}}{\log q},
\]
i.e.\ exactly the $\zeta$ law with the large-ordinate constant $4L^{3}/3$, divided by the
$\log q$ of the last table row. Equivalently both read
$(\text{window length})^{3}/6$ in their own units.

\paragraph{Which factors of $2$ cancel.} Two enter, in opposite directions, and they cancel:
\begin{enumerate}\itemsep1pt
\item \emph{Pair versus quartet.} Theorem~\ref{thm:ff} moves a reciprocal \emph{pair}; the
$\zeta$ counterfactual moves a \emph{quartet}, twice as many zeros --- a factor $2$ up.
\item \emph{$\sin^{2}$ averaging.} The $\zeta$ constant uses
$K(L,\gamma)\to L^{3}/3$, half of the $\int u^{2}=2L^{3}/3$ that the function-field centred
moment is the discrete analogue of --- a factor $2$ down.
\end{enumerate}

\paragraph{Checked on one curve.} For $y^{2}=x^{3}+x$ over $\Fq$, $q=5$, with
$a=10^{-6}$, the ratio of the exact function-field $\lmin$ to the $\zeta$-law prediction
$-\tfrac{4}{3}L^{3}\delta^{2}/\log q$ is
\[
1.406,\quad 1.195,\quad 1.096,\quad 1.047,\quad 1.024
\qquad\text{at } R=8,16,32,64,128,
\]
converging to $1$ like $1+3/R$ --- which is precisely
$\binom{R+2}{3}\big/\tfrac{R^{3}}{6}=(1+\tfrac1R)(1+\tfrac2R)$, the only discrepancy being
the $O(R^{2})$ term dropped in passing from the binomial to $R^{3}/6$. The dictionary is
therefore exact at leading order, with a fully accounted correction.

\section{The $\zeta$ side as a control}\label{sec:zeta}
\tier{T2; \src{src/stage3\_converge.py}, \src{src/stage3\_connes.py}}

We assemble the window form from Zhu's geometric side in three closed-form bases (two even,
one odd), moving everything to $u$-space by Parseval so no oscillatory quadrature is needed,
and summing the archimedean $m$-series analytically via digamma, Lerch $\Phi$ and Hurwitz
$\zeta$. The odd-sector pole term $-2P_jP_k$ --- the sign flips because
$F(-i/2)=-F(i/2)$ for odd $f$ --- is derived here.

\paragraph{Gate A.} Under RH the explicit formula gives
$Q(f)=\sum_\rho|F(\gamma_\rho)|^{2}$. The assembled geometric side matches the zeros side to
$1.63\times10^{-7}$ (even, vanishing at $\pm L$), $2.26\times10^{-7}$ (odd) and
$4.23\times10^{-3}$ (full cosine, which jumps at $\pm L$) at $K=800$ zeros, each decaying at
the rate that basis's smoothness predicts ($\gamma_K^{-3}$ versus $\gamma_K^{-1}$).
\emph{Control:} flipping the pole sign inflates the discrepancy to $5.35$ and $0.212$, so the
test discriminates.

\paragraph{Gate B.} $\lmin(N)$ decreases monotonically in two independent complete even
bases to $2.2702\times10^{-17}$ and $2.3357\times10^{-17}$ at $N=24$, landing on Zhu's
certified $[8.9\times10^{-18},\,2.27\times10^{-17}]$. Zhu normalises by $Q(f)/\|f\|^{2}$, the
same normalisation used here. We report the inequality, not the matching digits: the
sequence is still falling.

\paragraph{Connes's window.} At $2L=\log13$, $N=36$, $140$ digits:
$\lmin=8.977\times10^{-52}$, simple and even, with the ground state's Fourier zeros matching
$\gamma_1$ to $7.2\times10^{-48}$ and tracking roughly $N/2$ leading zeros. This reproduces
at smaller scale what Groskin~\cite{Groskin1} computes at $N=100$--$250$.

\paragraph{Simplicity and evenness.} \tier{T2, $N=12$, 40 digits} At
$L\in\{0.5,0.8,1.0,\tfrac12\log13\}$ the ground state is simple with relative gap
$2\times10^{4}$ to $4\times10^{5}$ and \emph{even} in every case, the odd sector sitting
$2$--$5$ orders above. The $L=0.8$ entry is consistent with Zhu's certificate; the other
three carry no external certificate and are T2 only.

\section{Detection onset}\label{sec:onset}
\tier{T2; \src{src/stage3e\_certify.py}, \src{data/stage3e\_certified\_onset.csv}}

\paragraph{Precedent.} Bombieri~\cite[\S13]{Bombieri} already ran this experiment in 2000:
the first $N\le160$ zeros of $\zeta$ together with ``a fictitious zero $\rho_0$ off the
critical line, together with their images by complex conjugation and reflection'' --- the
same quartet --- at $\rho_0=0.52+3.14i$, on windows $E=[-t,t]$, and found a critical value
$t_c$ beyond which the negative eigenvalue stays bounded away from $0$. He also separates
even and odd eigenfunctions. Our contribution is to resolve that critical window in
$\delta$ and in basis size, and to identify it with Zhu's bandwidth $T^{*}(L)=2\pi e^{2L}$.

\paragraph{The certification asymmetry.} A negative Rayleigh quotient from \emph{any} trial
function certifies indefiniteness. A finite basis that finds none certifies nothing. We
therefore grow the basis (low block $\cup$ resonant block $\{k:w_k\approx\gamma\}$) and
report, for each planted ordinate, the smallest window at which \emph{some} trial function
goes negative. The geometric side supplies the on-line part exactly; truncating the zeros
sum omits positive mass $\sim10^{-6}$ against a signal $\sim10^{-27}$ and cannot certify
anything near the onset.

\begin{center}
\begin{tabular}{@{}lcccc@{}}
\toprule
planted at & $\gamma_n$ & $L_{\mathrm{pred}}=\tfrac12\log\frac{\gamma}{2\pi}$ &
onset, small basis & onset, largest basis \\
\midrule
$\gamma_1$    & $14.13$  & $0.405$ & $1.3\times$ & $1.3\times$ \\
$\gamma_5$    & $32.94$  & $0.828$ & $1.0\times$ & $1.0\times$ \\
$\gamma_{10}$ & $49.77$  & $1.035$ & not found   & $1.0\times$ \\
$\gamma_{30}$ & $101.3$  & $1.390$ & not found   & $\mathbf{0.85\times}$ \\
$\gamma_{80}$ & $201.3$  & $1.733$ & not found   & $1.15\times$ \\
\bottomrule
\end{tabular}
\end{center}

The onset tracks $L_{\mathrm{pred}}$ within about $\pm30\%$ and moves \emph{earlier} as the
basis grows. Rows whose magnitude falls below the measured noise floor (the geometric
control is PSD by construction, so a negative value there \emph{is} the floor) are excluded
and flagged in the data file.

\paragraph{What detection actually needs --- and what it does not.} Two conditions are
visible in the data: the window must reach the ordinate,
$L\gtrsim L_{\mathrm{pred}}(\gamma_n)$, and the basis must \emph{contain modes resonant with}
$\gamma_n$, i.e.\ indices near $k^{*}\approx\gamma_n L/\pi$ --- which is what the resonant
block supplies. It does \emph{not} require a basis large enough to annihilate the first $n$
zeros. The dimension needed is far below the top mode index: $\gamma_{30}$ is detected with
dimension $46$ and top index $45$, and $\gamma_{80}$ with dimension $59$ and top index
$\approx134$, whereas annihilating the first $n$ zeros in a contiguous basis would need
$N\approx2n$, i.e.\ $60$ and $160$. The $N/2$ rule of \S\ref{sec:zeta} governs a different
question --- how many leading zeros a \emph{contiguous} basis can annihilate --- and is not
the detection criterion. Detection only needs some $f$ in the subspace with
$Q_{\mathrm{geom}}(f)+2F(\gamma)^{2}<4\delta^{2}\bigl(F'^{2}+FF''\bigr)$, which the sparse
``low block $\cup$ resonant block'' subspace achieves at a small fraction of the dimension a
contiguous basis would need.

\paragraph{Weak $\delta$-dependence, explained.} Detection is
$C(L)\delta^{2}>\lstar(L)$. By Zhu's Landau--Widom law $\lstar(L)$ falls
super-exponentially in $L$ while $C(L)\sim L^{3}$ grows polynomially, so a factor in
$\delta^{2}$ shifts the crossing only logarithmically. Measured: a tenfold change in
$\delta$ moves the onset from $L=0.5$ to $L=0.6$.

\section{The recorder split}\label{sec:recorder}
\tier{T2; \src{src/task2\_inertia\_reconcile.py}}

Write $\T=A-2M$. In the function field the pole part of $A$ is
$q^{|n|/2}+q^{-|n|/2}=u_iw_j+w_iu_j$, an exact rank-$2$ \emph{hyperbolic} plane whose scale
grows like $q^{R/2}$ against a fixed $2g\cdot I$; so $A$ is eventually indefinite under both
conventions tested, $M$ is indefinite, and $S=\T$ is PSD. This is structural.

On $\zeta$ windows it is not. With
$A=\mathrm{Pole}+\mathrm{Arch}-\log\pi\cdot I$, $2M=\mathrm{Prime}$ one finds $A$ with
negative index $1$ for $L\le0.8$ and $2$ at $L=\tfrac12\log13$, identical in all three
bases. But grouping $-\log\pi$ with the primes instead makes \textbf{both $A$ and $M$
positive definite} at $L\le0.8$. The reason is structural: for even $f$ the $\zeta$ pole
term is $h(i/2)+h(-i/2)=2F(i/2)^{2}\ge0$, rank one and \emph{definite}, not hyperbolic.
There is no hyperbolic plane on the $\zeta$ side to find. We therefore state the convention
($A=\mathrm{Pole}+\mathrm{Arch}-\log\pi\cdot I$) and claim Block F only for the function
field.

\section{Relation to prior work}\label{sec:prior}

\textbf{Hallouin--Perret~\cite{HP19}} is the function-field foundation: Proposition~5 is the
window Toeplitz matrix, Theorem~6 the first-singular-size and kernel generator,
Theorem~36(ii) step~(i), Lemmas~33--34 the rank rule and parity split used here.
\textbf{Connes--van Suijlekom~\cite{CvS}} supply the $C^{*}$-algebraic proof of step~(i) and
the Hurwitz step, cite~\cite{HP19} for the function-field connection, and state the
simplicity difficulty themselves; their Remark~2.3 is our degeneracy observation in words.
\textbf{Bombieri~\cite{Bombieri}} is the precedent for planted $\zeta$ experiments
(\S13) and proves the negative-eigenvalue count (Theorem~8); \textbf{Yoshida~\cite{Yoshida}}
first studied the Weil hermitian form on functions supported in $[-t,t]$.
\textbf{Zhu~\cite{Zhu}} supplies the certified $\lstar(0.8)$, the normalisation and the
bandwidth $T^{*}$. \textbf{Groskin~\cite{Groskin1,Groskin2}} supplies the large-scale
Galerkin computations and the archimedean tail order.
\textbf{Connes--Consani--Moscovici~\cite{CCM}}, \textbf{Connes~\cite{ConnesLetter}},
\textbf{Connes--Consani~\cite{CC21}} and \textbf{Suzuki~\cite{Suzuki}} are the surrounding
program. A fuller account, with verbatim quotations and the verification method, is in
\texttt{PRIOR\_ART.md}.

\section{Limitations}\label{sec:limits}

\begin{enumerate}\itemsep2pt
\item No rigorous lower bound on $\lstar(L)$ of our own. Theorem~\ref{thm:zeta} is a
\emph{relative} estimate; it becomes an unconditional lower bound only where Zhu's certificate
reaches ($L\le0.8$). Replacing $\lstar(L)$ by $0$ beyond that range would assume the window
positivity at issue, and is not asserted (Retraction~18).
\item Theorem~\ref{thm:zetasharp} is conditional on all other zeros being on the line.
Nonnegativity of the remaining-zero form $\widetilde Q=Q_\zeta-2F(\gamma)^{2}$ is not supplied
by positivity of $Q_\zeta$, by positivity of $Q_0$, or by a zero-free strip.
\item A zero-location restriction alone does not discharge the hypothesis of
Theorem~\ref{thm:zetasharp}. For the audited source version and the transfer attempted in
\src{extensions/openai-\allowbreak quasi-rh-\allowbreak audit-2026-10-07/}, no usable estimate
was obtained that
replaces the remaining-zero positivity hypothesis: the restriction bounds how far a zero may
sit from the critical line, not the aggregate window-weighted contribution of all such zeros,
and the per-zero bound used there is uniform in the ordinate and therefore not summable as it
stands. This is a limitation of that particular estimate, not a general obstruction.
\item Interval certification covers the zeros-side assembly only: \texttt{mpmath}'s interval
context has no digamma and no Lerch transcendent, so the geometric side could not be
interval-evaluated. What is certified instead is the LDL$^{\mathsf T}$ witness identity
$v^{\mathsf T}Zv=D_j$ and a measured noise floor.
\item Onset non-detections are limited by the subspaces actually searched: the largest had
dimension $59$ with top mode index $\approx134$. A non-detection bounds that search, not the
method.
\item $C(L)$ is measured at finite $\delta$ and finite basis; the saturation
$C/(4L^{3}/3)\to1$ is observed, not proved.
\item The $\zeta$-side numbers depend on the planted ordinate, the basis and the truncation.
\end{enumerate}

\section{Traps this project documents}\label{sec:traps}

The retraction log (Appendix~\ref{app:retr}) is part of the result. The reusable lessons:
\textbf{(a)} truncating the zeros sum biases the form \emph{downward} and manufactures false
detections; \textbf{(b)} moving a zero and adding a zero are two different counterfactuals,
both legitimate --- Bombieri studies the second --- and the sharp law needs the first;
\textbf{(c)} \texttt{float64} returns spurious negative eigenvalues of size $10^{-16}$ where
the true value is exactly $0$; \textbf{(d)} interval libraries have gaps, and
\texttt{mpmath}'s \texttt{iv.sincpi} is itself broken (it calls a non-existent
\texttt{ctx.sinpi}); \textbf{(e)} a non-detection certifies nothing; \textbf{(f)} keeping
known-good regression values was the only thing that caught three sector-rename bugs.

\appendix
\section{Retractions}\label{app:retr}

Eighteen corrections were logged during this work; the full list with evidence is in
\texttt{REPORT.md}. Those a reader must see:

\begin{center}
\begin{longtable}{@{}cp{0.84\textwidth}@{}}
\toprule
\# & Retraction \\
\midrule
\endhead
1 & Frobenius radius computed on $P$ rather than its reverse ($0.8$ instead of
$\le1.3\times10^{-51}$). \\
3 & A tolerance-based sign test misclassified exact zeros; replaced by an error-based test
plus exact rational certificates. \\
4 & \textbf{H2 as first stated was refuted by my own adversarial pass}; restated with a
distinctness hypothesis. The restatement is Theorem~\ref{thm:HP6}. \\
6 & \textbf{``Step (iv) is free'' retracted} --- it had been scored on zero angles, which is
the wrong object. At the function level it carries content
(Proposition~\ref{prop:hurwitz}). \\
8--9 & A 40-digit run discarded; three sector-rename bugs caught only by regression values. \\
11 & \textbf{``Low-lying zeros only'' retracted as structural.} The height sweep had used a
truncated on-line side, and non-detections were basis-limited. Redone,
$\gamma_{30}\approx101$ is detected at $0.85\,L_{\mathrm{pred}}$. \emph{Amended:} the
added-quartet counterfactual is legitimate in its own right --- it is Bombieri's --- and was
not the error. \\
12 & $\gamma_{80}$ ``detections'' below their own noise floor discarded; later detected at
$1.15\,L_{\mathrm{pred}}$ with a positive control. \\
15 & \textbf{``Block F holds on the $\zeta$ side'' retracted} (\S\ref{sec:recorder}). \\
16 & \textbf{Three claimed contributions are published theorems}: step~(i) is
\cite[Thm.~36(ii)]{HP19}, the first-singular-size result is \cite[Thm.~6]{HP19}, the window
matrix is \cite[Prop.~5]{HP19}; the degeneracy observation is \cite[Rem.~2.3]{CvS}; the
negative-index count is \cite[Thm.~8]{Bombieri}; the planted experiment is
\cite[\S13]{Bombieri}. An earlier abstract-level novelty check had missed all of these. \\
17 & \textbf{The $N\gtrsim2n$ detection rule retracted} (\S\ref{sec:onset}). It contradicted
both the stated limitation and the $\gamma_{80}$ table row. The $\gamma_{80}$ detection used
dimension $59$, top index ${\approx}134$, against $2n=160$; $\gamma_{30}$ used dimension
$46$, top index $45$, against $2n=60$. The rule had been imported from \S\ref{sec:zeta},
where $N/2$ governs a stronger question --- how many leading zeros a \emph{contiguous} basis
annihilates. The table row and the limitation were correct; only the rule was wrong. \\
18 & \textbf{The unsupported replacement of the unperturbed baseline by zero outside the
certified range has been withdrawn.} Theorem~\ref{thm:zeta} previously added that for
$L>0.8$ it ``holds with $\lstar(L)$ replaced by $0$, using no positivity input at all''.
Substituting $0$ for $\lstar(L)$ in a lower bound requires $\lstar(L)\ge0$, which is the
finite-window Weil positivity at issue; at $\delta=0$ the substituted statement asserts
$\lmin(Q_0)\ge0$, i.e.\ the very input it disclaimed. The preceding relative perturbation
estimate is retained, and is now stated as such. The same step was the last link of the chain
in \S B.4 of \src{notes/proofs.md} and is corrected there. No numerical value, no other
theorem and no table is affected. \\
\bottomrule
\end{longtable}
\end{center}

\section{Reproduction}\label{app:repro}

Code and data: \url{https://github.com/innerlightr-wq/weil-window-control}.
Python~3.12, \texttt{mpmath}~1.3.0, \texttt{numpy}~1.26.4, \texttt{matplotlib}~3.8+;
no \texttt{gmpy2} (installing it speeds the high-precision runs substantially).

\begin{center}
\begin{tabular}{@{}lll@{}}
\toprule
Result & Script & Data \\
\midrule
\S\ref{sec:planted} genuine & \src{src/stage1.py} & \src{stage1\_spectra.csv} \\
\S\ref{sec:planted} planted & \src{src/stage2.py} & \src{stage2\_spectra.csv} \\
\S\ref{sec:planted} kernel & \src{src/function\_level.py} & \src{stage2\_function\_level.csv} \\
Thm~\ref{thm:ff}, \ref{thm:zeta}, \ref{thm:zetasharp} & \src{src/proofs\_check.py} & \src{stage3\_sharp\_bound.csv} \\
$C$ vs Cauchy--Schwarz & \src{src/task34\_extra.py} & \src{stage3\_C\_vs\_CS.csv} \\
$K(L,\gamma)$, dictionary & \src{src/task\_final\_checks.py} & \src{stage3\_exact\_constant\_K.csv} \\
Gates A, B & \src{src/stage3\_converge.py} & \src{stage3\_gateA/B\_*.csv} \\
Connes window & \src{src/stage3\_connes.py} & \src{stage3b\_connes\_zeros\_N36.csv} \\
Detection onset & \src{src/stage3e\_certify.py} & \src{stage3e\_certified\_onset.csv} \\
Recorder split & \src{src/task2\_inertia\_reconcile.py} & \src{stage3d\_inertia\_reconciled.csv} \\
\bottomrule
\end{tabular}
\end{center}

\section*{Acknowledgments}
\addcontentsline{toc}{section}{Acknowledgments}

AI assistance was used for algebra checking, numerical verification, adversarial review,
literature organization, and manuscript drafting. All conceptual development, research
direction, interpretation, and final responsibility for the content belong to the author.

\begin{thebibliography}{99}

\bibitem{Bombieri} E. Bombieri,
\emph{Remarks on Weil's quadratic functional in the theory of prime numbers, I},
Atti Accad. Naz. Lincei Cl. Sci. Fis. Mat. Natur. Rend. Lincei (9) Mat. Appl. \textbf{11}
(2000), no.~3, 183--233.

\bibitem{CC21} A. Connes and C. Consani,
\emph{Weil positivity and trace formula the archimedean place},
Selecta Math. (N.S.) \textbf{27} (2021), no.~4, art.~77.
\href{https://doi.org/10.1007/s00029-021-00689-4}{doi:10.1007/s00029-021-00689-4};
arXiv:2006.13771.

\bibitem{CCM} A. Connes, C. Consani and H. Moscovici,
\emph{Zeta spectral triples}, arXiv:2511.22755.

\bibitem{ConnesLetter} A. Connes,
\emph{The Riemann Hypothesis: Past, Present and a Letter Through Time}, arXiv:2602.04022.

\bibitem{CvS} A. Connes and W. D. van Suijlekom,
\emph{Quadratic forms, real zeros and echoes of the spectral action}, arXiv:2511.23257.

\bibitem{Groskin1} A. Groskin,
\emph{High-precision approximation of Riemann zeros via the truncated Weil form},
arXiv:2605.20224. \href{https://doi.org/10.5281/zenodo.19546514}{doi:10.5281/zenodo.19546514}.

\bibitem{Groskin2} A. Groskin,
\emph{A finite Guinand--Weil dictionary and archimedean tail order for the truncated Weil
quadratic form}, arXiv:2607.02828.

\bibitem{HP19} E. Hallouin and M. Perret,
\emph{A unified viewpoint for upper bounds for the number of points of curves over finite
fields via Euclidean geometry and semi-definite symmetric Toeplitz matrices},
Trans. Amer. Math. Soc. \textbf{372} (2019), no.~8, 5409--5451.
\href{https://doi.org/10.1090/tran/7813}{doi:10.1090/tran/7813}.

\bibitem{Suzuki} M. Suzuki,
\emph{Weil's quadratic form via the screw function}, arXiv:2606.09096.

\bibitem{Yoshida} H. Yoshida,
\emph{On Hermitian forms attached to zeta functions},
in: N. Kurokawa and T. Sunada (eds.), \emph{Zeta Functions in Geometry},
Adv. Stud. Pure Math. \textbf{21}, Math. Soc. Japan, Kinokuniya, Tokyo, 1992, 281--325.

\bibitem{Zhu} X. Zhu,
\emph{Weil positivity in compact windows: a finite reduction, certified two-sided bounds,
and a Landau--Widom decay law}, arXiv:2608.24827.

\bibitem{PartitionNote} E. De Jes\'us,
\emph{The Partition Potential and the Laguerre Tower}, Zenodo, 2026.
\href{https://doi.org/10.5281/zenodo.21108392}{doi:10.5281/zenodo.21108392}.

\bibitem{EnforcerNote} E. De Jes\'us,
\emph{Riemann Hypothesis: The Weil Enforcer in Audit Coordinates}, Zenodo, 2026.
\href{https://doi.org/10.5281/zenodo.21115524}{doi:10.5281/zenodo.21115524}.

\end{thebibliography}

\end{document}
